% !PRINTED 231
\weberpage{231}

%% [running head, not body text -- commented out 2026-08-07] Page furniture (running
%% head, folio, signature) is excluded from this edition by standing policy.
%% \S.\,74. \hfill Hauptgleichung fünften Grades \hfill 231

%<<CONT
\weq{2}{
\begin{array}{rcl@{\qquad}rcl}
5\,B_{0,0} &=& 4\,a_4^2 - 10\,a_3\,a_5, & 5\,B_{0,1} &=& 3\,a_3\,a_4 \\[1ex]
5\,B_{1,1} &=& 6\,a_3^2, & 5\,B_{0,2} &=& 0 \\[1ex]
5\,B_{2,2} &=& -20\,a_0\,a_4, & 5\,B_{0,3} &=& -25\,a_0\,a_5 \\[1ex]
5\,B_{3,3} &=& 0 & 5\,B_{1,2} &=& -25\,a_0\,a_5 \\[1ex]
5\,B_{1,3} &=& -20\,a_0\,a_4, & 5\,B_{2,3} &=& -15\,a_0\,a_3.
\end{array}
}
Wenn man hieraus die Determinante berechnet, so erhält
man ohne Schwierigkeit nach \S.\,73, (12) für die Discriminante
der Hauptgleichung fünften Grades
\[
\text{(3)} \qquad D = a_0^2\bigl(5^5\,a_0^2\,a_5^4 + 2^8\,a_0\,a_4^5 + 2250\,a_0\,a_3^2\,a_4\,a_5^2
- 1600\,a_0\,a_3\,a_4^3\,a_5 + 108\,a_3^5\,a_5 - 27\,a_3^4\,a_4^2\bigr)\,.
\]
Da nun hier $B_{3,3} = 0$ ist, so haben wir ein Werthsystem
der $t$, wofür die Bezoutiante verschwindet, nämlich
\[
t_0 = 0, \quad t_1 = 0, \quad t_2 = 0\,.
\]
Wir nehmen nun unsere Tschirnhausen-Substitution in folgender Form [\S.\,68, (12)]
\[
\text{(4)} \qquad y = t_3\,F_0 + t_2\,(\alpha_2\,F_1 + \alpha_1\,F_2 + \alpha_0\,F_3)
\]
und wollen über die $\alpha_0, \alpha_1, \alpha_2$ so verfügen, dass die Gleichung
fünften Grades für $y$ unabhängig von $t_3, t_2$ zu einer Hauptgleichung wird, dass also $S(y^2)$ identisch für alle $t_2, t_3$ verschwindet.

\indent Da $S(F_0^2) = B_{3,3}$ hier verschwindet, so geschieht dieser
Forderung nach \S.\,73, (11) genüge, wenn wir $\alpha_0, \alpha_1, \alpha_2$ so bestimmen, dass
\[
\text{(5)} \qquad \alpha_0^2\,B_{0,0} + \alpha_1^2\,B_{1,1} + \alpha_2^2\,B_{2,2}
+ 2\,\alpha_0\,\alpha_1\,B_{0,1} + 2\,\alpha_0\,\alpha_2\,B_{0,2} + 2\,\alpha_1\,\alpha_2\,B_{1,2} = 0\,,
\]
\[
\text{(6)} \qquad \alpha_0\,B_{0,3} + \alpha_1\,B_{1,3} + \alpha_2\,B_{2,3} = 0\,.
\]
Diese Gleichungen führen durch Elimination von einer der
drei Grössen $\alpha$ auf eine quadratische Gleichung für das Verhältniss der beiden anderen. An sich ist es nicht nothwendig,
irgend einen besonderen Fall auszuschliessen, auch nicht den,
dass $B_{0,3}, B_{1,3}, B_{2,3}$ alle drei verschwinden; man hätte dann nur
für die $\alpha$ irgend eine Lösung der Gleichung (5) zu nehmen.

\indent Um aber nicht weitläufig zu sein, wollen wir annehmen,
dass $B_{0,3}$ von Null verschieden sei, weil die hierdurch ausgeschlossenen Fälle, in denen eine Wurzel der gegebenen Gleichung Null oder unendlich wird, hier eigentlich kein Interesse
bieten, weil sie auf eine Gleichung niedrigeren Grades zurückkommen. Die Behandlung bleibt übrigens ganz dieselbe, wenn
%CONT>>
